\section{Dual heights and the structural coefficient}\label{sec:heights}

Throughout the proof, \(b>0\) denotes the spin inverse temperature.
The height spacing is \(2\pi\). On an original, or \emph{primary}, dual
edge the weight of the difference \(2\pi j\) is
\[
 p_b(j)=e^{-b}I_j(b),\qquad
 \sum_{j\in\Z}p_b(j)e^{tj}=\exp\{b(\cosh t-1)\}.
\]
Thus \(p_b\) is the law of the difference of two independent
Poisson variables of mean \(b/2\). An unoriented edge is counted once.
For a retained primary graph, \(A\) is its unit-conductance Laplacian:
\[
 \ip{g}{Ag}=\sum_{\{x,y\}}(g_x-g_y)^2.
\]
This normalization converges to \(\int|\nabla g|^2\).
Pairings on a finite graph use counting measure.
A free connected height component is taken modulo common translations
by \(2\pi\Z\); one may choose a pinned representative when evaluating
translation-invariant tests. When confining observations control the
constant, all its translation copies are instead summed.

\subsection{Duality and finite comparisons}

The current-to-height correspondence below is the standard planar
XY duality; see \cite[Section~2]{vEL} and
\cite{vanEngelenburgLisDuality2025}. We give the finite free-box
identity to fix its boundary and winding conventions.

\begin{lemma}[Free-boundary duality]\label{lem:duality}
In a finite free spin rectangle, the two-point spin correlation is
the ratio of a twisted height partition sum to the centered height
partition sum, both with weights \(p_b\) and exterior face height zero.
The twist is an integral connection with circulations \(+2\pi,-2\pi\)
around the two source vertices and zero circulation elsewhere.
For interior sources it can vanish near the outer boundary.
After cuts, the same description holds on each retained chart, without
identifying heights through the cut.
\end{lemma}

\begin{proof}
Expand each factor \(e^{b\cos(\theta_x-\theta_y)}\) in its Fourier
series. Integrating the spins imposes zero divergence on the integer
edge current in the denominator. Inserting
\(e^{i(\theta_0-\theta_x)}\) changes its divergence to the two unit
sources. The insertion with the opposite sign has the same sum by
current reversal, so taking the real part makes no change.
The factors \(e^b\) cancel between the two partition sums.

Rotate currents onto dual edges. In a planar rectangle, every
divergence-free integral current is a difference of integer face
heights with exterior height zero. This follows by summing the
rotated current along dual paths: path independence is exactly the
zero-divergence condition at the enclosed primal vertices. Subtract
a fixed unit current on a path between the sources in the numerator.
The remaining current is divergence free and gives the same height
variables. The subtracted current is the claimed connection.
Changing its path is an integral change of height variables.
Equivalently, the numerator has multivalued heights with the stated
monodromies and single-valued branches on simply connected charts.
Releasing identifications across a cut gives the asserted cut version.
\end{proof}

The finite Gaussian inequalities of \Bcite{st:finite-gaussian},
based on the lattice inequality of Regev and
Stephens-Davidowitz \cite{RSD}, apply to this nonquadratic weight
through a series approximation. We record the extension, including
the forms needed after conditioning.
The general use of Gaussian representations for this transfer
has precedent in \cite[Lemmas~9.6--9.7 and Appendices~C--D]{AHPS2022};
the Bessel mixture representation there is credited to Aran Raoufi.
The series construction below verifies the particular affine
constraints and observation forms used here. It does not require
the bare Bessel potential to satisfy a super-Gaussian condition
at every value of \(b\); subdivision is used in such
extensions, as also in \cite[Theorem~8 and Section~1.3.3]{LammersDichotomy2026}.

\begin{lemma}[Finite comparison rules]\label{lem:comparisons}
On a fixed primary graph, allow additional nonnegative quadratic
forms in real linear observations, centered equalities, and
independent Gaussian observation noises.
\begin{enumerate}[label=\textup{(\roman*)},leftmargin=*]
 \item Increasing centered precision decreases centered covariance
 and every centered real moment generating function. Covariance on
 a translate, including after a real linear tilt, dominates centered
 covariance. In particular
 \[
  \log\E e^X\ge \tfrac12\Var(X)
 \]
 for a centered inversion-symmetric law and a real linear test \(X\).
 \item A shifted-to-centered partition ratio is at most one and
 decreases when the same positive precision is added to the affine
 and centered coordinates. Shifts may be nonzero equality values,
 an integral connection, or centers of observation penalties.
 The monotonicity also holds with a fixed real linear statistic in
 the numerator exponent, but the bound by one need not hold for
 this tilted ratio. Equality constraints may be inserted by
 infinite penalties whenever the affine constraint is nonempty.
 \item Centered characteristic coefficients are nonnegative and
 increase with precision. Their logarithms have Hessian at every
 phase at least the Hessian at zero. A phase on a free component is
 allowed if its coefficient total is an integer. In the invariant
 finite-graph situations used below these coefficients are positive.
\end{enumerate}
Released constant modes are handled modulo translations or by
confining observations. Real tests used when releasing a constant
annihilate that constant.
\end{lemma}

\begin{proof}
Replace every primary edge by \(N\) bonds in series, each with
normalized integer increment weight
\[
 q_N^{j^2},\qquad q_N=b/(2N).
\]
For all sufficiently large \(N\) this is a genuine lattice Gaussian.
For every fixed real \(t\), uniformly on compact sets of \(t\),
\[
 \frac{\sum_jq_N^{j^2}e^{tj}}{\sum_jq_N^{j^2}}
 =1+\frac bN(\cosh t-1)+O_t(N^{-2}).
\]
The sum of the \(N\) increments therefore converges to \(p_b\),
with convergence of exponential moments at every fixed real
argument. Matching chain endpoints and imposing the cycle and
period constraints gives a quadratic lattice law on the enlarged
graph. The numerator is an affine lattice parallel to the
denominator lattice.

All partition ratios just listed converge at fixed primary graph.
For domination, use a spanning forest to express vertex differences
as sums of chain increments; the remaining normalized edge weights
are bounded by one. The forest variables have uniformly bounded
exponential moments. If a constant is confined, sum its translations
against the confining Gaussian. This also dominates fixed real
tilts and polynomial factors. Normalizers of the independent chain
laws agree in the two sectors.

On the enlarged lattice, the covariance and partition comparisons
are precisely \Bcite{st:finite-gaussian}. They hold on a proper
affine subspace by restricting to its parallel span and completing
the square in the remaining fixed coordinates. For clarity, with
quadratic precision \(T+tQ\), \(Q\ge0\), differentiation of the log
shifted-to-centered ratio gives
\[
 -\frac12\tr Q\bigl(
       \Cov_{\mathrm{shift}}-\Cov_0
       +m_{\mathrm{shift}}\otimes m_{\mathrm{shift}}\bigr)\le0.
\]
After a fixed real tilt, completing the square gives the same
formula with the tilted covariance and mean. This proves the
tilted version of (ii), not merely its untilted version.
Infinite centered penalties implement the required equalities.
Poisson summation on the enlarged lattice proves (iii).
Taking the fixed-graph limit proves the stated rules.
Vanishing auxiliary masses justify released translation modes;
alternatively one can use integer gradient coordinates.
\end{proof}

The bounds that remain uniform in subdivision are not the bare
Gaussian bounds of the enlarged network. They are bounds on whole
chain increments.

\begin{lemma}[Reference-cost bounds]\label{lem:reference-cost}
Suppose a primary test has a flow representation
\[
 X=\sum_e w_e\nabla_e h,\qquad
 \mathcal R(w)=\sum_e w_e^2.
\]
In a centered law with any additional centered precision,
\begin{equation}
 \Var(X)\le C_b\mathcal R(w),\qquad
 \E|X|^r\le C_{r,b}\mathcal R(w)^{r/2}.
 \label{eq:reference-moments}
\end{equation}
For each fixed \(t\),
\(\E\exp(tX/\sqrt{\mathcal R(w)})\) is bounded uniformly in the graph
and in sufficiently fine subdivision. The constants can be uniform
for \(b\) in a compact subset of \((0,\infty)\).
These statements include minimization over flows with the prescribed
divergence and grounded vertices.
\end{lemma}

\begin{proof}
Keep independent full-chain increments as variables and release
their endpoint, curl, and period constraints. Lemma~\ref{lem:comparisons}
bounds the original centered transform by this product transform.
For one subdivision bond and bounded \(t\), pairing opposite
increments in its moment generating function gives
\[
 \log\frac{\sum_jq_N^{j^2}e^{2\pi tj}}
               {\sum_jq_N^{j^2}}
 \le C_{b,R}t^2/N\qquad (|t|\le R).
\]
Summing over \(N\) bonds gives \(C_{b,R}t^2\) for a whole chain.
For a flow normalized by \(\sqrt{\mathcal R(w)}\), every coefficient
has absolute value at most one. Independence therefore bounds its
log transform at a fixed real argument by a constant. This proves
the normalized exponential and moment bounds, and rescaling proves
\eqref{eq:reference-moments}. The same inequalities on the original
law follow by the fixed-graph limit.
In this argument a centered primary observation is a real linear
function of the enlarged coordinates, so its squared penalty is
an additional positive form. Cycle, period, and endpoint conditions
are likewise limits of positive equality penalties. They can all
be removed before applying the independent whole-chain bound.
No uniform estimate on the inverse microscopic Gaussian precision
is asserted or used.
\end{proof}

We call \(\mathcal R\), or its minimum over flows, the
\emph{reference cost}. In particular it always refers to the primary
unit-edge graph. For a zero-total profile \(f\), this minimum is
\(\ip{f}{A^+f}\), by the usual orthogonal decomposition of flows into
a minimizing gradient flow and a divergence-free flow.

There are two useful further consequences. First, if a phase
coefficient \(\phi(s)>0\) is varied by a zero-total test \(Y\) of
variance at most \(V\), the logarithmic Hessian bound gives
\begin{equation}
 \sqrt{-\log\phi(s+tY)}
 \le \sqrt{-\log\phi(s)}+|t|\sqrt{V/2}.
 \label{eq:phase-transfer}
\end{equation}
Indeed \(f=\log\phi\le0\), \(f''\ge-V\), imply
\(|f'|^2\le-2Vf\), which integrates to the displayed inequality.
This proves the same transfer bound as \Bcite{st:coefficient-transfer},
now with centered variance and hence reference cost.
An integer-total phase on a fixed free graph is equivalent to a
zero-total phase after subtracting an integer point charge.
Equation~\eqref{eq:phase-transfer} starting at the zero phase proves
its strict positivity.

Second, centered restoration comparisons can be localized.
Represent a removed edge by an independent chain and restore it
by the centered equality matching its endpoint difference.
Adding exterior bonds or centered guard pins increases the
probability of these equalities. The centered restoration ratio is
therefore between its value in a free patch and its value with a
separating guard pinned to zero. At that pin the interior no longer
depends on the exterior, provided no observation crosses the guard.
If all endpoints of the restored bonds, and the support of any
restored form, lie in \(H\), pinning every primary vertex of \(H\)
makes restoration a fixed scalar, also for a chain. Bayes' rule and
centered precision comparison bound the free-to-guard multiplicative
gap by
\begin{equation}
 \frac{\Prob(h_p=0\mid h_H=0)}{\Prob(h_p=0)}.
 \label{eq:centered-restoration-gap}
\end{equation}
These assertions also follow by interpolating positive forms before
taking equality limits. Independent-chain scalar normalizers are
the same in each member of the comparison.

\subsection{Free Gaussian limits}

For \(B_m=\{0,\ldots,m-1\}^2\), put \(g_m(x)=x_1/m\) and
\[
 a_m(b)=\Var_{B_m,b}\ip{h}{A_{B_m}g_m}.
\]
No bare quadratic parameter is attached to this definition.

\begin{theorem}[The free coefficient]\label{thm:free-limits}
There exists a finite \(a=a(b)\ge0\) such that \(a_m(b)\to a\).
Let \(B_n\), at mesh \(n^{-1}\), be any sequence of free axis
rectangles with endpoints converging to those of a nondegenerate
rectangle \(B\). For every finite collection of smooth shifts
defined near \(\overline B\),
\begin{equation}
 \log\E\exp\{\ip{h}{A_{B_n}g(\cdot/n)}\}
 \longrightarrow \frac a2\int_B|\nabla g|^2.
 \label{eq:free-gradient-limit}
\end{equation}
Joint limits are Gaussian, with covariance
\(a\int_B\nabla g\cdot\nabla g'\), and all joint moments converge.
For zero-total profiles \(f_n\) whose scaled densities converge in
\(L^2(B)\) to a bounded piecewise continuous \(f\) of integral zero,
\begin{equation}
 \log\E e^{\ip{h}{f_n}}
 \longrightarrow
 \frac a2\ip{f}{(-\Delta_{B,N})^{-1}f}_{L^2(B)}.
 \label{eq:free-smear-limit}
\end{equation}
Face smears, after subtracting their totals from a volume smear,
are allowed as well. These limits may be degenerate when \(a=0\).
\end{theorem}

\begin{proof}
We give the changes to the cutting and folding proof of
\Bcite{st:free-rectangle-limit}; the observation Laplacian here is
the primary \(A\), not the subdivided quadratic precision.

Cut an \(M\)-square into \(m\)-squares and strips of width less than
\(m\). In its affine observation, the seam and strip flows have
reference cost \(O(m^{-1}+m/M)\). Cutting increases covariance, and
the retained cell observations become independent. Consequently
\begin{equation}
 \sqrt{a_M}\le\sqrt{a_m}
             +C_b(m^{-1}+m/M)^{1/2}.
 \label{eq:cut-coefficient}
\end{equation}
First let \(M\to\infty\), and then let \(m\) follow a liminf
subsequence. This proves existence and finiteness of \(a\).
Square symmetry makes the limiting covariance of the two affine
observations scalar.

Cutting rectangles into squares gives the upper variance bound
\(a\int_B|\nabla g|^2\) for affine \(g\).
For the reverse bound, at each fixed discrete rectangle tile an
arbitrarily large square by copies of that rectangle, choosing its
side divisible by both rectangle side counts. Apply the same seam
estimate, first sending that square side to infinity. If
\(m_1(n),m_2(n)\) are the rectangle side counts and \(V_n(p)\) is its
variance for \(g=p\cdot x/n\), this gives
\[
 \sqrt a\,|p|
 \le\frac{n}{\sqrt{m_1(n)m_2(n)}}\sqrt{V_n(p)}
       +\frac{C_{b,p}}{\sqrt{\min(m_1(n),m_2(n))}}.
\]
Thus every affine shift attains the upper bound along every
admissible rectangle sequence.

Cut into cells of small fixed macroscopic diameter and approximate
each gradient there by a constant. Seams have vanishing reference
cost, and the interior error is controlled by the squared
\(L^2\)-modulus of continuity. This proves
\begin{equation}
 \limsup_n\Var\ip{h}{A_{B_n}g}
       \le a\int_B|\nabla g|^2
 \label{eq:partial-energy-upper}
\end{equation}
for smooth or continuous piecewise smooth \(g\), and for partial
gradient observations restricted to polygonal pieces.
Call equality in this upper bound \emph{saturation}.
Saturated shifts form a linear space: on a convergent covariance
subsequence the upper-bound defect is positive semidefinite, and
its diagonal vanishes for saturated shifts, so every entry vanishes.
Saturation passes to an enclosed rectangle by cutting its complement
into rectangles and comparing the sum of their upper bounds with
the saturated outer variance. It also passes through gradient-energy
approximation, by Lemma~\ref{lem:reference-cost}.

To obtain nonaffine saturated shifts, condition on the whole trace
of a site-axis or diagonal reflection line and pin one of its
vertices. The subdivided quadratic graph is attractive:
its density satisfies the lattice condition edge by edge.
Its conditioned law is associated, and association passes to the
fixed-primary-graph limit. If a saturated seed is odd about the
line, the two partial tests on its open sides are reflected
negatives. If the off-line coefficients on one side are
nonnegative after a vanishing-cost transport to the line,
conditional association gives nonpositive conditional cross
covariance. The two conditional means are negatives, so the
unconditional cross covariance is nonpositive as well, up to the
vanishing transport error. The saturated sum and the separate
upper bounds \eqref{eq:partial-energy-upper} force each partial
test to saturate. The corresponding folded shift differs from
that partial test only along a vanishing-cost seam.

Here the coefficient verification can be made on unit edges.
For the axis seed \(g=x_1\), the only negative coefficients of the
right-side partial observation occur next to the cut. Move them
to the conditioned line along bounded-length paths. There are
\(O(n)\) coefficients of size \(O(n^{-1})\), with bounded overlap,
so the cost is \(O(n^{-1})\).
The hinges \((x_1-t)_+\), and then all smooth univariate functions,
are therefore saturated.
For a diagonal seed, put \(Y=y+t\) and take
\(g=F(x)-F(Y)\) on a symmetric enclosing square
\(p\le x,Y\le q\). Initially assume \(F'''<0\) and
\(F'(p),F'(q)>0\). In \(x>Y\), away from the diagonal strip,
the interior coefficient is nonnegative because the symmetric
second differences of \(F\) decrease with their center.
The coefficients on \(x=q\) and \(Y=p\) have positive leading
terms \(F'(q)/n\) and \(F'(p)/n\); their \(O(n^{-2})\)
remainders are harmless. Move the exceptional diagonal-strip
coefficients to the line. Their total path cost is again
\(O(n^{-1})\). The fold criterion proves saturation of
\begin{equation}
 (F(x)-F(y+t))\,1_{\{x>y+t\}}.
 \label{eq:diagonal-fold}
\end{equation}
Every smooth \(F\) is a difference of two functions with the strict
properties just used: add a sufficiently large cubic with negative
third derivative and then a sufficiently large positive linear
function. Hence \eqref{eq:diagonal-fold} holds for arbitrary \(F\).
Coherently round a reflection line by \(O(n^{-1})\).
The changed fold has squared gradient cost \(O(n^{-1})\):
only an \(O(n^{-1})\)-width strip changes its side, and away from
that strip the smooth gradient changes by \(O(n^{-1})\).
Restriction from the enclosing symmetric rectangle consequently
gives the claim on every admissible target rectangle sequence.

Finally choose a smooth family \(F_t\), compactly supported in the
relevant \(t\)-range, with \(F_t'(x)=-g_{xy}(x,x-t)\).
The mixed derivative of the integral of
\eqref{eq:diagonal-fold} is \(g_{xy}\). Its difference from \(g\)
is therefore a sum of univariate functions. Riemann sums converge
in gradient \(L^2\), including across the moving diagonals.
Linearity and energy approximation prove saturation for every
smooth \(g\).

The lower Laplace bound now follows from
Lemma~\ref{lem:comparisons}. For the upper bound, cut into cells of
side comparable to a fixed \(r\). Remove seam errors by Holder's
inequality and the reference-cost exponential bound. A cell test
has reference-cost square root \(O(r)\); hence, for fixed \(p>1\),
\[
 \log\E e^{pX}
   =\tfrac12p^2\Var X+O_p(r^3).
\]
The remainder follows from
\(\E[|X|^3e^{p|X|}]=O_p(r^3)\), which uses normalized exponential
moments, not a bare Gaussian upper tail.
There are \(O(r^{-2})\) cells. First take the microscopic limit,
then \(r\downarrow0\), then \(p\downarrow1\).
This proves \eqref{eq:free-gradient-limit}; applying it to real
linear combinations proves the joint assertion and moment
convergence.

Nearest-neighbor Poincare bounds the reference cost of a density
error by a constant times its scaled \(L^2\)-norm squared.
Approximate \(f\) by finite sums of Neumann cosines, use their
smooth inverse-Laplacian shifts, and then remove the approximation.
For such a shift \(n^2A_{B_n}g\) converges in density \(L^2\) to
\(-\Delta g\); the boundary coefficients are \(O(n^{-2})\) and
occupy \(O(n)\) sites. Using the actual rectangle endpoints before
letting them converge handles arbitrary admissible sequences.
This proves \eqref{eq:free-smear-limit}.
A face probability can be moved to a strip of width \(r\) at
reference cost \(O(r+n^{-1})\), by parallel normal flows.
Strip approximation proves the face extension.
\end{proof}

\subsection{Soft joining and the coefficient gap}

Write \(u_m\) for uniform probability on \(B_m\), and
\[
 w_m(s)=\E_{B_m,b}e^{i\ip{h}{s}},
 \qquad \sum_xs_x\in\Z.
\]
The definition is independent of the representative modulo
\(2\pi\Z\), and \(0<w_m(s)\le1\).

\begin{theorem}[Uniform phase and forbidden range]
\label{thm:coefficient-gap}
For every \(b>0\),
\[
 a(b)=0\quad\hbox{or}\quad a(b)\ge8\pi.
\]
If \(a(b)>0\), the fractional mean
\(\ip{h}{u_m}\bmod2\pi\) converges to uniform measure on the circle,
independently of any fixed finite collection of the Gaussian
observations of Theorem~\ref{thm:free-limits}.
More generally, for a fixed nonzero integer \(l\) and any zero-total
profiles \(t_m\) of bounded reference cost,
\begin{equation}
 w_m(lu_m+t_m)\longrightarrow0.
 \label{eq:uniform-phase}
\end{equation}
\end{theorem}

\begin{proof}
We first isolate the replacement for the quadratic charge-space
joining used in \Bcite{st:fractional-mean} and
\Bcite{st:critical-face-ratio}. Tile a rectangle by free cells,
keeping their lattice translations, and independently assign every
seam edge a difference of law \(2\pi p_b\).
Let \(\mathcal V\) be the neutral Euclidean space of cell totals.
Choose a linear flow right inverse \(P\) on the neighboring-cell
graph, so \(\operatorname{div}P(d)=d\) for \(d\in\mathcal V\).
Define \(H\in\mathcal V\) by
\begin{equation}
 \ip{H}{d}
 =\sum_{\langle i,j\rangle}P(d)_{ij}
   \left(\ip{h_i}{\rho_{ij}}-\ip{h_j}{\rho_{ji}}
                            -\overline\eta_{ij}\right).
 \label{eq:soft-mismatch}
\end{equation}
Here \(\rho_{ij}\) is the uniform seam-face probability and
\(\overline\eta_{ij}\) is the mean of the independent seam
differences, with consistent orientation. Translating the cells
by constants \(c_i\) changes \(H\) by the orthogonal projection
of \(c\) onto \(\mathcal V\).
Weight this law by \(e^{-|H|^2/2}\), summing relative translations
and taking the common translation modulo \(2\pi\).
Adding exact seam matching makes \(H=0\) and recovers the joined
rectangle. The soft law thus has larger internal variances and
smaller invariant characteristic coefficients than the joined law.
These comparisons with translation-summed laws can be justified
without biasing the initial free-cell laws. First add the auxiliary
precision \(\eta\sum_i\overline h_i^2\), where \(\overline h_i\)
is the cell mean. For a fixed cell shape with representative mean
\(\mu\), its translation sum is
\(\sum_{k\in\Z}e^{-\eta(\mu+2\pi k)^2/2}\).
After division by its Gaussian integral scalar, this tends
uniformly in \(\mu\) to one and is uniformly bounded, by Poisson
summation. Thus the original internal law tends to the product
of the free-cell laws, also in moments of invariant tests.
After adding \(|H|^2\), divide out the common translation scalar.
The relative translation sum is a periodized Gaussian on the
projected lattice, uniformly bounded in its shape-dependent
shift. Dominated convergence again applies with the internal
moments. Apply precision comparison at \(\eta>0\) and then let
\(\eta\downarrow0\). In particular the soft variance of the sum
of two internal affine tests is at most their independent
free-cell variance.

For a phase whose cell restrictions are \(s_i\), of totals
\(\xi_i\), Poisson summation in the relative translations gives,
up to a common scalar, the numerator
\begin{equation}
 \sum_{\substack{n_i\in\Z\\\sum n_i=\sum\xi_i}}
 e^{-|d|^2/2}
 \prod_i w_i\left(s_i+\sum_{j\sim i}P(d)_{ij}\rho_{ij}\right)
 \chi_m(P(d)),\qquad d=n-\xi.
 \label{eq:soft-fourier}
\end{equation}
The denominator uses \(s=0\). The signs of \(P\) are chosen so
the displayed cell totals are \(n_i\).
Indeed the reciprocal vectors of the projected translation
lattice are exactly the neutral integer vectors; the phase
shifts them to \(n-\xi\). Fourier transforming the Gaussian
soft weight supplies \(e^{-|d|^2/2}\); the remaining factors
are the cell and independent-seam characteristic functions.
The latter factor \(\chi_m\) is positive and at most one.
For fixed cell fluxes it tends to one, since each seam average
has variance \(O(m^{-1})\).
For two cells one may use a scalar mismatch penalty, giving the
same formula with integer sectors \(l\) and weights \(e^{-l^2/2}\).
All these calculations can first be performed on the subdivided
quadratic graph and then passed to the fixed primary graph.

We prove the uniform phase. Let \(\rho\) be the right-face
probability, and set
\[
 X=\ip{h}{A(x_1/m)},\qquad
 Y=\ip{h}{\rho-u_m}.
\]
On the unit square graph \(X\) is exactly the right-face average
minus the left-face average. Theorem~\ref{thm:free-limits} gives
joint Gaussian limits, \(\Var X\to a\), and, by reflection,
\(\Cov(X,Y)\to a/2\).
Suppose, for some nonzero integer \(l\), that
\(w_m(lu_m)\) stays bounded away from zero on a subsequence.
The reference-cost exponential bounds give a further subsequence
on which
\[
 F_m(t,v)=\E e^{il\ip{h}{u_m}+itX+ivY}
\]
converges locally with all derivatives to an entire function
\(F\). On real arguments it is positive: this follows from
\eqref{eq:phase-transfer} and the assumed positive value at the
origin.

Here is the replica calculation responsible for its form.
For two independent copies of a centered subdivided quadratic
height, put \(S=h_1+h_2\), \(D=h_1-h_2\).
Conditioned on their common parity class modulo \(4\pi\Z\),
they are independent and identically distributed.
Writing \(V_{ij}=\E(S_iS_j\mid\text{parity})\), we have
\begin{equation}
 \E V_{ij}=2\Cov(h_i,h_j),\qquad
 \Cov(V_{ij},V_{rs})=2\kappa(h_i,h_j,h_r,h_s).
 \label{eq:replica}
\end{equation}
To verify the second identity, put \(C_{ij}=\E h_ih_j\) and
\(M_{ijrs}=\E h_ih_jh_rh_s\).
Conditional independence and expansion in the two original copies
give
\[
 \E V_{ij}V_{rs}
 =\E S_iS_jD_rD_s
 =2M_{ijrs}+2C_{ij}C_{rs}
             -2C_{ir}C_{js}-2C_{is}C_{jr}.
\]
Subtracting \(\E V_{ij}\E V_{rs}=4C_{ij}C_{rs}\) gives exactly
twice the fourth joint cumulant, without symmetrizing the indices.
In particular
\(\Var(V_{XX})=2\kappa(X,X,X,X)\) and
\(\Var(V_{XY})=2\kappa(X,X,Y,Y)\).
After first taking the fixed-graph chain limit and then the
Gaussian limit, both vanish.

For completeness, write \(F_{m,N}\) for the transform before the
chain limit, and put
\(\theta(h)=l\ip{h}{u_m}+tX+vY\).
Inversion symmetry makes \(F_{m,N}\) real on real arguments.
For the observation indices \(i,j\in\{X,Y\}\), differentiating
in their respective phase variables gives
\[
 \E S_iS_j e^{i\theta(D)}
 =-2\bigl(F_{m,N}F_{m,N,ij}
                     -F_{m,N,i}F_{m,N,j}\bigr).
\]
By conditional independence the left side equals
\(\E[V_{ij}\E(e^{i\theta(D)}\mid\text{parity})]\).
Replacing \(V_{ij}\) by \(2C_{ij}\) has absolute error at most
\(\sqrt{\Var(V_{ij})}\), and
\(\E e^{i\theta(D)}=F_{m,N}^2\).
Fixed-graph moment convergence passes this bound through
\(N\to\infty\); Gaussian moment convergence then removes the
error as \(m\to\infty\). Thus
\[
 -2FF_{tt}+2F_t^2=2aF^2,\qquad
 -2FF_{tv}+2F_tF_v=aF^2.
\]
Since \(F>0\) on real arguments, division by \(-2F^2\) yields
\[
 (\log F)_{tt}=-a,\qquad(\log F)_{tv}=-a/2.
\]
Reflection gives \(F_t(0,0)=0\); hence
\begin{equation}
 F(t,v)=F(0,v)\exp(-at^2/2-atv/2).
 \label{eq:phase-transform}
\end{equation}

Now softly join two horizontal \(m\)-squares and test the sum
of their internal affine observations. Its variance is at most
twice the free-cell variance, and at least its variance in the
fully joined rectangle. The joined rectangle's full affine
test differs from this sum by the seam test, whose reference
cost is \(O(m^{-1})\): there are \(m\) seam coefficients of
size \(m^{-1}\). Its limiting variance is consequently \(2a\),
so the soft-law variance also tends to \(2a\).

In \eqref{eq:soft-fourier}, every sector function has logarithmic
second derivative at least minus twice the free-cell variance.
The logarithmic second derivative of a positive sum is the
weighted mean of these derivatives plus the variance of the
sector logarithmic first derivatives. Thus the negative second
derivative of the normalized sum, the soft-law variance, is at
most twice the free variance minus that nonnegative sector
variance. Differentiation of the sector sum is justified by
the following summable bounds. Write \(\psi_l(t)\) for the
product of the two cell characteristic coefficients in sector
\(l\), and \(V=2\Var(X)\) for the independent-cell variance.
Then \(0<\psi_l\le1\) and
\((\log\psi_l)''\ge-V\), so the derivative estimate underlying
\eqref{eq:phase-transfer} gives
\[
 \psi_l|(\log\psi_l)'|^2
       \le 2V(-\psi_l\log\psi_l)\le 2V/e.
\]
Multiplication by \(e^{-l^2/2}\chi_m(l)\le e^{-l^2/2}\)
gives a summable bound, uniform in \(m\), on the
logarithmic-derivative square terms. The raw second derivatives
are bounded by the independent-cell second moment \(V\).
The analogous first-derivative bound follows by
Cauchy--Schwarz. Hence termwise differentiation and the
mixture-Hessian identity are valid even for sectors whose
coefficients become small.
The sector zero has weight one before normalization. The sector
\(l\) stays bounded away from zero: transfer from \(lu_m\) to
the seam face costs a bounded amount, and \(\chi_m\to1\).
By \eqref{eq:phase-transform} at \(v=l\), reflection and charge
reversal give its logarithmic first derivative tending to
\(-al\), up to the irrelevant common orientation. The zero
sector derivative is zero. If \(a>0\), this contradicts the
vanishing sector variance. Therefore \(w_m(lu_m)\to0\).
Equation~\eqref{eq:phase-transfer} proves
\eqref{eq:uniform-phase}; Fourier approximation on the circle
and the Gaussian limits prove the joint uniform-phase assertion.

It remains to obtain the sharp numerical gap.
Let \(\overline\rho_m\) be one quarter of the sum of the four
face probabilities and put \(z_m=w_m(\overline\rho_m)>0\).
Tile a \(Gm\)-square into \(G^2\) cells, with \(G\) fixed.
In the outer-face phase, the cell totals \(\xi_i\) are supported
at the boundary, sum to one, and satisfy
\(\sum_i\xi_i^2\le C/G\).
The denominator of \eqref{eq:soft-fourier} tends to one:
its zero sector is one; every other sector has a cell of
nonzero integer total and tends to zero by
\eqref{eq:uniform-phase}; Gaussian sector weights permit
dominated convergence at fixed \(G\).

Retain numerator sectors \(n=e_{i_0}\), where \(i_0\) is at
least \(G/4\) from the boundary. The radial-flow construction
of \Bcite{st:radial-cost} supplies outward face fluxes \(H_{ie}\)
with divergence \(e_{i_0}\), constant outer flux \(1/(4G)\),
and, on putting
\[
 D_{ij}=(H_{i,+j}-H_{i,-j})/2,\qquad
 S_{ij}=H_{i,+j}+H_{i,-j},
\]
the bounds
\begin{equation}
 \sum_{i,j}D_{ij}^2\le\frac1{2\pi}\log G+C,\qquad
 \sum_{i,j}S_{ij}^2\le C.
 \label{eq:radial-cost}
\end{equation}
For clarity, its continuum field is
\[
 V(x)=\frac{x-x_0}{2\pi|x-x_0|^2}+E(x),\qquad
 |E|\le C/G,\quad|\nabla E|\le C/G^2.
\]
The correction fixes the constant outer normal flux; it is
obtained by the zero-Neumann problem after even reflection.
Integrating \(V\) across unit-cell faces gives \(H\).
The squared radial field integrates to
\((2\pi)^{-1}\log G+O(1)\); face-sum derivative errors are
summable and give the second bound in \eqref{eq:radial-cost}.
These constants are uniform for the retained central cells.
The internal flows have divergences \(e_{i_0}-\xi\).
These vectors are linearly independent, since their central
coordinates are distinct and \(\xi\) is supported on boundary
cells. Prescribe \(P\) on them by these flows and extend
linearly. Thus one soft precision works for all retained
sectors.

In each cell, expand the logarithm of its characteristic
coefficient around the symmetric four-face profile of the
same total, either zero or \(\overline\rho_m\).
The linear term vanishes by square symmetry. The logarithmic
Hessian lower bound therefore gives a product lower bound
\[
 z_m\exp\left\{-\frac12\sum_i\Var\ip{h}{Y_i}\right\},
\]
where \(Y_i\) is the face-profile difference.
Coordinate reflections split \(Y_i\) into two mutually
orthogonal odd face differences and an even remainder.
The odd differences have limiting variance \(a\); the even
remainder has reference cost bounded by
\(C\sum_jS_{ij}^2\). Consequently, as also calculated in
\Bcite{st:face-energy-bound},
\[
 \limsup_m\sum_i\Var\ip{h}{Y_i}
       \le a\sum_{i,j}D_{ij}^2+C_b\sum_{i,j}S_{ij}^2.
\]
The soft penalty is bounded below uniformly on these
sectors, since \(|e_{i_0}-\xi|^2\) is bounded, and the seam
factor tends to one at fixed \(G\).
There are at least \(cG^2\) retained sectors. Characteristic
monotonicity from the soft law to the joined law yields
\begin{equation}
 \liminf_{m\to\infty}\frac{z_{Gm}}{z_m}
       \ge c_bG^{\,2-a/(4\pi)},
 \label{eq:face-ratio}
\end{equation}
with \(c_b>0\) independent of sufficiently large fixed \(G\).
If \(0<a<8\pi\), choose one integer \(G\) making the right side
larger than two. Then \(z_{Gm}\ge(3/2)z_m\) for every
sufficiently large \(m\). Iterating along \(m,Gm,G^2m,\ldots\)
contradicts \(0<z_m\le1\). This proves the gap.
\end{proof}

\subsection{An upper bound for spin correlations}

\begin{proposition}[Spin upper exponent]\label{prop:spin-upper}
For the prescribed free-box correlation,
\begin{equation}
 \limsup_{n\to\infty}\frac{\log C_b(n)}{\log n}
       \le-\frac{2\pi}{a(b)}.
 \label{eq:spin-upper}
\end{equation}
When \(a(b)=0\), the assertion means that the limsup is
\(-\infty\), or equivalently that every fixed polynomial
upper exponent is eventually available.
\end{proposition}

\begin{proof}
In a sufficiently large free spin box, put disjoint circular
dyadic annuli of radii \(r,2r\) around each source, through
radii of order \(n\). On the fixed annulus \(1<|z|<2\),
choose a smooth tangential divergence-free form \(p\),
compactly supported there and approximating \(d\arg z\)
in \(L^2\). A radial cutoff provides such forms.
Discretize \(p\) at scale \(r\) by taking differences of a
sampled stream function on faces. Its pairing with twisted
differences is deterministic, by circulation, and tends to
\[
 D=\int p\cdot\nabla\arg z.
\]
Use the opposite sign at the negative source.

Insert in the twisted numerator the real exponential of
\(\lambda\) times the sum of these deterministic pairings.
Release endpoint constraints between a grid of small
rectangular cells at every annular scale, retaining
independent seam increments; release all remaining exterior
constraints as well. Lemma~\ref{lem:comparisons}(ii), with
its real tilt, bounds this normalized tilted numerator by
the resulting product of transforms.
Each retained cell is simply connected and avoids the
sources, so its integral connection disappears by an
integer gauge change.
Here the tilt is on the affine edge differences, the same
coordinates as the edge weights. If the connection on a chart
is \(\alpha=\nabla\kappa\), the reindexing \(h'=h+\kappa\)
replaces \(\nabla h+\alpha\) by \(\nabla h'\) in both the
weight and the tilt. It therefore introduces no residual
real-tilt scalar.

Take cells of side comparable to \(tr\), lying within the
annulus, with \(t>0\) small and fixed. There are only
\(O_t(r)\) leftover nonzero edge coefficients, each
\(O(r^{-1})\); their independent-chain log moment is
\(o_r(1)\). On the squares, approximate the coefficients by
constant gradients. Theorem~\ref{thm:free-limits}, Holder's
inequality, and the reference-cost bound on the coefficient
error give log transform at most
\[
 \frac{\lambda^2a}{2}\int|p|^2+o_r(1)+o_t(1).
\]
The errors are taken first at fixed \(t,\lambda,p\); rounding
changes only vanishing reference-cost terms.

There are \(\log n/\log2+O(1)\) annuli per source.
At any fixed \(n\) all cuts are internal to a sufficiently
large box, so we may first take its prescribed free-box
limit. Averaging the convergent estimates along the
increasing annular scales gives
\[
 \limsup_n\frac{\log C_b(n)}{\log n}
 \le\frac{2}{\log2}
       \left(-\lambda D+\frac{\lambda^2a}{2}\int|p|^2\right),
\]
after \(t\downarrow0\).
Now let \(p\to d\arg z\); both \(D\) and
\(\int|p|^2\) tend to \(2\pi\log2\).
For \(a>0\), the optimal choice \(\lambda=1/a\) yields
\(-2\pi/a\). The factor two is from the two sources.
For \(a=0\), each fixed \(\lambda>0\) is allowed, giving
\(-4\pi\lambda\); letting this fixed parameter be arbitrarily
large proves the stated interpretation.
\end{proof}
